\scanpage{52}
\mspage{51}
2 - 7 - \hfill n° 392

Pour tout $x\in K$, soit $\Delta(x)$ l'élément $(x_{\underline p})\in\prod_{\underline p\in P}\hat K_{\underline p}$ tel que $x_{\underline p}=x$ pour tout $\underline p\in P$ ; comme $x_{\underline p}\in\hat A_{\underline p}$ sauf pour un nombre fini de valeurs de $\underline p$, on a $\Delta(x)\in\underline A$ ; on définit donc ainsi un homomorphisme $\Delta:K\longrightarrow\underline A$, qui est \uline{injectif} si $P\ne\varnothing$, (c'est-à-dire si $A$ n'est pas un corps) : les éléments de $\Delta(K)$ sont dits \uline{adèles restreints principaux}, et il est clair que $\Delta(A)\subset\underline A_0$.

PROPOSITION 3.- \uline{L'anneau} $\underline A_0$, (resp. $\underline A$), \uline{s'identifie au complété de} $A$, (resp. $K$), \uline{pour la topologie d'anneau dont un système fondamental de voisinages de} $0$ \uline{est formé de tous les idéaux entiers} $\ne(0)$ \uline{de} $A$.

Il est immédiat que la topologie considérée sur $A$, (ou $K$), est séparée. Compte tenu du n° 3, l'assertion relative à $\underline A_0$ résulte du chap. III, § 2, n° 13, prop. 17. Cela montre donc que $\Delta(A)$ est dense dans $\underline A_0$ ; pour voir de même que $\Delta(K)$ est dense dans $\underline A$, on remarque que pour tout $\underline x=(x_{\underline p})\in\underline A$, il n'y a qu'un nombre fini de $\underline p\in P$ tels que $v_{\underline p}(x_{\underline p})<0$ ; en vertu du § 1, n° 5, prop. 9, il y a donc un $s\in K$ tel que $sx_{\underline p}\in\hat A_{\underline p}$ pour tout $\underline p\in P$, autrement dit $\Delta(s)\underline x\in\underline A_0$, et comme la multiplication par $\Delta(s)$ est un homéomorphisme de $\underline A$ sur lui-même, il suffit d'appliquer le fait que $\Delta(A)$ est dense dans $\underline A_0$ pour en déduire que $\Delta(K)$ est dense dans $\underline A$.

On pourrait naturellement aussi prouver que $\Delta(K)$ est dense dans $\underline A$ en utilisant la prop. 2.

Considérons maintenant le groupe multiplicatif $\mathrm{SL}(n,\underline A)$, formé des matrices $\underline U\in M_n(\underline A)$ telles que $\det(\underline U)=1$ ; si on munit $M_n(\underline A)=\underline A^{n^2}$ de la topologie produit, elle induit sur $\mathrm{SL}(n,\underline A)$ une topologie \uline{compatible avec la structure de groupe} de $\mathrm{SL}(n,\underline A)$. En effet, il suffit de voir que l'application $\underline U\longrightarrow\underline U^{-1}$ est continue dans $\mathrm{SL}(n,\underline A)$ ; mais comme $\underline U$ est unimodulaire, on sait, (\uline{Alg.}, chap. III, 3e éd., §\quad, n°\quad), que les éléments de $\underline U^{-1}$ sont des \uline{mineurs} de $\underline U$, donc des polynômes en les éléments de $\underline U$, ce qui prouve notre assertion. Si on identifie $K$ à un sous-anneau
